\section{Заключение}
\label{sec:Chapter5} \index{Chapter5}

Существуют и другие схожие решения, в частности кольцевой буфер, используемый
в eBPF~\cite{ebpf_ringbuf}, в котором этап резервирования реализован с применением спинлоков.
Такой подход обеспечивает корректность конкурентной записи, однако в условиях
виртуализации использование блокирующих примитивов синхронизации может
приводить к существенным задержкам.
В случае переключения контекста внутри гостевой операционной системы удерживаемая
блокировка препятствует прогрессу остальных потоков-писателей. Своё влияние на
непредсказуемость задержек также оказывает планирование виртуальных процессоров
на стороне гипервизора, что может существенно увеличить время удержания блокировки.

Преимуществом предложенного решения перед механизмом, используемым в Linux eBPF,
является возможность отобразить кольцевой буфер непосредственно в виртуальное
адресное пространство пользовательского процесса. Благодаря этому запись
диагностических сообщений может выполняться без обращения к драйверу,
что снижает накладные расходы на системные вызовы и переключение контекста.

В рамках данной работы был разработан альтернативный механизм передачи
диагностических данных, основанный на неблокирующем MPSC-кольцевом буфере,
размещённом в памяти гостевой операционной системы. Для синхронизации
конкурентных записей используется вышеописанный lock-free алгоритм с
разделением операций резервирования, копирования и фиксации записи
(commit), что позволяет полностью отказаться от блокирующих примитивов
синхронизации; такой подход может быть полезен и при одновременной работе
нескольких виртуальных машин на одном хосте. Предложенное решение обеспечивает независимый прогресс
потоков и корректную работу даже при условии приостановки выполнения
отдельных потоков-писателей. Реализация интегрирована в виртуализированную
среду, представленную эмулятором Stratovirt, UEFI EDK2 и гостевой
операционной системой Windows. Данный механизм позволяет обеспечить
более предсказуемое поведение при высокой интенсивности диагностических
событий.
